\documentclass{hogent-article}

% Invoegen bibliografiebestand
\addbibresource{voorstel.bib}

\studyprogramme{Professionele bachelor toegepaste informatica}
\course{Bachelorproef}
\assignmenttype{Onderzoeksvoorstel}
\academicyear{2026-2027}


\title{Hoe kan een uitlegbare risicoscore voor phishingmails worden opgebouwd uit structurele e-mailkenmerken?}
\author{Pieter Swillens}
\email{pieter.swillens@student.hogent.be}
\supervisor[Co-promotor]{Y. Nijs (Yonicon BV, \href{mailto:youri.nijs@yonicon.be}{youri.nijs@yonicon.be})}


% Binnen welke specialisatierichting uit 3TI situeert dit onderzoek zich?
% Kies uit deze lijst:
%
% - Mobile \& Enterprise development
% - AI \& Data Engineering
% - Functional \& Business Analysis
% - System \& Network Administrator
% - Mainframe Expert
% - Als het onderzoek niet past binnen een van deze domeinen specifieer je deze
%   zelf
%
\specialisation{Cybersecurity}
\keywords{Phishingdetectie, E-mailbeveiliging, Risicoscore}

\begin{document}
\begin{abstract}
Phishing via e-mail blijft volgens \textcite{Veit2025} een van de grootste beveiligingsrisico's voor zowel bedrijven als alledaagse internetgebruikers. Ook in België is de omvang aanzienlijk: fraudeurs maakten in 2024 ongeveer 49 miljoen euro buit \autocite{Febelfin2025}. De opkomst van GenAI maakt het bovendien eenvoudiger om overtuigende en gepersonaliseerde berichten op te stellen \autocite{SchmittFlechais2024}. Sensibilisering alleen volstaat niet, want gebruikers blijven ook na phishingtraining kwetsbaar \autocite{Marshall2024}. En wanneer ook technische beveiliging, zoals de filters van de mailclient, faalt, ligt de beoordeling van de legitimiteit bij de gebruiker. Op dat moment kan de gebruiker ondersteuning gebruiken bij het beoordelen van de e-mail.

Deze bachelorproef onderzoekt hoe een interpreteerbare risicoscore voor phishingmails kan worden opgebouwd uit structurele e-mailkenmerken, zodat phishing met een lage false-positive-rate wordt gedetecteerd en elke score herleidbaar is tot een controleerbaar bewijs in de e-mail. Het beoogde resultaat is een proof of concept dat een e-mail analyseert, er een risicoscore aan toekent en de achterliggende redenen in gewone taal weergeeft. Omdat uitleg een belangrijk onderdeel vormt, wordt gekozen voor een interpreteerbaar model, namelijk logistische regressie \autocite{Rudin2019}. Uit elke e-mail worden kandidaat-kenmerken berekend. Via L1-regularisatie, een embedded methode voor kenmerkselectie \autocite{Guyon2003}, wordt bepaald welke kenmerken in het uiteindelijke model blijven. 

Het model wordt getraind en geëvalueerd op een combinatie van Phishing Pot \autocite{PhishingPot} en de legitieme e-mail partities van het SpamAssassin Public Corpus \footfullcite{SpamAssassinCorpus}. De prestaties worden gemeten op een testset, met de \emph{false-positive-rate} (\emph{FPR}) als belangrijkste metriek. Verwacht wordt dat het model een \emph{FPR} van hoogstens 3\% en een precision van minstens 95\% behaalt, en dat elke risicoscore herleidbaar is tot een concreet, controleerbaar bewijs in de e-mail. Zo wordt nagegaan in hoeverre een uitlegbare risicoscore technisch haalbaar is.
\end{abstract}

\tableofcontents
%---------- Inleiding ---------------------------------------------------------
\section{Inleiding}%
\label{sec:inleiding}
Cybercriminaliteit die gericht is op het verkrijgen van geld of persoonsgegevens is geen recent verschijnsel. Reeds in de jaren negentig slaagden aanvallers erin om via phishing de authenticatiegegevens van America Online (AOL)-gebruikers te bemachtigen \autocite{Goel2018}. Sindsdien heeft phishing zich verder ontwikkeld en is het uitgegroeid tot een veelvoorkomende vorm van cybercriminaliteit. De blijvende relevantie van deze vorm van criminaliteit blijkt onder meer uit recente cijfers over de omvang van phishing in België.

In 2025 stuurden burgers bijna 10 miljoen verdachte berichten door naar \url{verdacht@safeonweb.be}\footnote{Centrum voor Cybersecurity België, \url{https://ccb.belgium.be/nl}, geraadpleegd op 8 oktober 2026.}. De omvang van phishing is trouwens niet uitsluitend zichtbaar in het aantal gemelde berichten, maar ook in de financiële schade die ermee gepaard gaat. Volgens \textcite{Febelfin2025} maakten fraudeurs in 2024 ongeveer 49 miljoen euro buit.

Hoewel phishing via verschillende kanalen verloopt, blijft misleiding via e-mail volgens \textcite{Veit2025} een van de grootste beveiligingsrisico's voor zowel bedrijven als alledaagse internetgebruikers. Wanneer zo'n phishingmail de filters van de mailclient passeert, is het de verantwoordelijkheid van de gebruiker om in te schatten of de mail te vertrouwen is of niet. Deze beoordeling wordt verder bemoeilijkt door de opkomst van generatieve artificiële intelligentie, die aanvallers in staat stelt om overtuigende en gepersonaliseerde phishingberichten te creëren in een handomdraai \autocite{SchmittFlechais2024}.

Hoewel sensibilisering gebruikers kan helpen om phishingmails te herkennen, volstaat ze niet om het risico op succesvolle aanvallen volledig weg te nemen. Onderzoek toont aan dat gebruikers ook na phishing\-training kwetsbaar blijven voor dergelijke aanvallen \autocite{Marshall2024}. 
Daarom is aanvullende ondersteuning wenselijk op het moment dat de gebruiker een e-mail moet beoordelen. 

Deze bachelorproef onderzoekt daarom de volgende centrale onderzoeksvraag: \emph{``Hoe kunnen relevante kenmerken uit e-mails worden geïdentificeerd en gecombineerd om phishingmails automatisch te detecteren en hun risiconiveau te bepalen?''}

Naast de scriptie levert het onderzoek een algoritme op dat een e-mail analyseert, er een risicoscore aan toekent en de redenen om tot die score te komen op een begrijpelijke en toegankelijke manier aan de gebruiker toont. Om de centrale onderzoeksvraag te beantwoorden, worden de volgende deelvragen behandeld:

\begin{enumerate}
  \item \emph{Welke kenmerken zijn het meest relevant om phishing te detecteren?}
  \item \emph{Hoe kunnen deze kenmerken automatisch uit een e-mail worden gehaald?}
  \item \emph{Hoe kunnen deze kenmerken worden gecombineerd tot een risicoscore?}
\end{enumerate}

Het onderzoek wordt als geslaagd beschouwd wanneer het algoritme, geëvalueerd op een vooraf afgebakende dataset van phishingmails en legitieme e-mails, phishing betrouwbaar van legitieme e-mail onderscheidt, gemeten aan de hand van evaluatiecriteria die in de literatuurstudie worden onderbouwd.


%---------- Stand van zaken ---------------------------------------------------
\section{Literatuurstudie}%
\label{sec:literatuurstudie}

\subsection{Kenmerken voor de detectie van phishingmails}
\label{subsec:Kenmerken voor de detectie van phishingmails}
NOG TE DOEN: UITLEGGEN WAT EEN DETECTIEKENMERK IS!!!

\input{inhoud/literatuurstudie/1.kenmerken-voor-detectie-van-phishingmails/1.kenmerken-van-de-afzender-en-mailheaders.tex}

\input{inhoud/literatuurstudie/1.kenmerken-voor-detectie-van-phishingmails/2.kenmerken-van-inhoud-en-taalgebruik.tex}

\subsection{Extractie van detectiekenmerken uit e-mails}
\label{subsec:extractie-van-detectiekenmerken}

Zoals reeds besproken in sectie~\ref{subsubsec:Afzender en e-mailheaders} volgt een e-mail de standaard uit RFC~5322 \autocite{Resnick2008}. Een e-mail bestaat dus uit headers en een body, die meerdere onderdelen kan bevatten, zoals platte tekst, HTML en bijlagen. Een mailclient vertaalt dit onderliggende bericht naar een weergave die de gebruiker kan lezen en interpreteren. Echter kan daarbij informatie verloren gaan. Zo toont een mailclient vaak enkel de weergavenaam van de afzender en niet het volledige e-mailadres. Ook de zichtbare tekst van een hyperlink hoeft niet overeen te komen met de URL waarnaar de link daadwerkelijk verwijst \autocite{Veit2025}. De gebruiker ziet daardoor niet de werkelijke bestemming van de link, al kan deze wel uit de onderliggende HTML-code worden afgeleid. Volgens \textcite{Chen2020} kunnen aanvallers deze verschillen misbruiken om toch een legitiem ogende mail te fabriceren.

Een toepassing die een e-mail wil beoordelen, heeft dus het onderliggende bericht nodig en niet enkel de weergave ervan. Hieronder worden de mogelijke routes besproken:

\begin{enumerate}
    \item \textbf{\texttt{.eml}-bestanden.} Een \texttt{.eml}-bestand bevat een volledige e-mail in het formaat volgens RFC~5322 \parencite{Resnick2008}, inclusief de headers die de ontvangende server heeft toegevoegd. Dit formaat komt overeen met dat van de datasets uit sectie~\ref{subsubsec:datasets}.
    \item \textbf{Gmail API.} Mits toestemming van de gebruiker kan een toepassing een bericht rechtstreeks opvragen bij de mailprovider \footfullcite{GmailMessagesGet}. De e-mail kan zowel als origineel bericht of als een door Gmail gestructureerde weergave worden opgevraagd.
\end{enumerate}

De verkregen gegevens zijn niet altijd onmiddellijk bruikbaar voor kenmerkextractie, omdat de headers en de berichtinhoud verschillende tekensets en coderingen kunnen gebruiken \autocite{Koide2025}. Zo beschrijft RFC~2047 \autocite{Moore1996} hoe niet-ASCII-tekst in headers wordt weergegeven met behulp van zogenaamde \emph{encoded-words}, die moeten worden gedecodeerd om de oorspronkelijke tekst te achterhalen. Ook de wijze van aanlevering kan een codering toevoegen: de Gmail API levert het volledige bericht bij \texttt{format=raw} bijvoorbeeld als base64url-gecodeerde tekenreeks \footfullcite{GmailMessages}. Waar nodig worden de gegevens daarom eerst gedecodeerd voordat kenmerken eruit kunnen worden afgeleid.
\subsection{Methoden om phishingmails te detecteren}
\label{subsec:methoden-detectie}

Een afzonderlijk detectiekenmerk volstaat zelden om een e-mail als phishing te classificeren, aangezien dezelfde signalen ook bij legitieme berichten kunnen voorkomen. Een detectiesysteem moet daarom meerdere kenmerken combineren tot één beoordeling en bepalen hoe zwaar elk kenmerk doorweegt in de eindscore. Dit onderdeel bespreekt twee manieren om dat te doen: \emph{rule-based} detectie en detectie op basis van \emph{machine learning}.

\input{inhoud/literatuurstudie/3.methoden-voor-detectie-van-phishingmails/1.rulebased-methoden.tex}
\input{inhoud/literatuurstudie/3.methoden-voor-detectie-van-phishingmails/2.machinelearning-methoden.tex}

\subsection{Uitlegbaarheid en gebruikersondersteuning}
\label{subsec:uitlegbaarheid-en-gebruikersondersteuning}

Het doel van deze scriptie is de gebruiker ondersteunen in de beoordeling van potentiële phishingmails. Een cruciaal onderdeel van die ondersteuning is een uitleg over waarom een mail een bepaald risiconiveau krijgt. Die \emph{uitlegbaarheid} wordt in dit onderdeel besproken.

\input{inhoud/literatuurstudie/4.uitlegbaarheid-en-gebruikersondersteuning/1.vormen-en-eigenschappen-van-uitlegbaarheid.tex}
\input{inhoud/literatuurstudie/4.uitlegbaarheid-en-gebruikersondersteuning/2.uitleg-voor-gebruikers.tex}
\subsection{Evaluatiecriteria}
\label{subsec:evaluatiecriteria}

De prestaties van een detectiesysteem hangen sterk af van twee keuzes: de data waarop het wordt getest en de maten waarmee de uitkomst wordt beoordeeld \autocite{Salloum2022}. Dit onderdeel bespreekt daarom eerst de beschikbare datasets en hun beperkingen, en vervolgens de gangbare evaluatiematen.

\input{inhoud/literatuurstudie/5.evaluatiecriteria/1.datasets.tex}
\input{inhoud/literatuurstudie/5.evaluatiecriteria/2.evaluatie-risicoscore.tex}


%---------- Methodologie ------------------------------------------------------
\section{Methodologie}
\label{sec:methodologie}

Om de onderzoeksvraag te kunnen beantwoorden, wordt er een proof of concept ontwikkeld.

\subsection{Requirements}
\label{subsec:requirements}
Deelvraag 1 vraagt aan welke eisen een risicoscore moet voldoen om gebruikers te ondersteunen bij het beoordelen van e-mails. Het effect op gebruikers valt buiten de scope van dit onderzoek. De eisen worden afgeleid uit sectie~\ref{subsubsec:uitleg-voor-gebruikers}.

\begin{description}
  \item[R1: weinig \emph{false positives}.] Gebruikers wennen aan foutieve waarschuwingen die ze herhaaldelijk zien \autocite{Anderson2015}. False positives moeten daarom zo veel mogelijk vermeden worden.
  \item[R2: aanwijzingen bij de bron.] Waarschuwingen werken beter dicht bij de verdachte link \autocite{Petelka2019}. De uitleg moet gebruikers dus kunnen wijzen op de plaats in de e-mail waar de aanwijzing staat.
  \item[R3: begrijpelijke taal.] Veel gebruikers begrijpen technische waarschuwingen niet \autocite{Egelman2008}. De uitleg moet dus in gewone taal en zonder jargon zijn geschreven.
  \item[R4: geen overdreven zekerheid.] Een overtuigende uitleg kan tot \emph{overreliance} leiden \autocite{Bucinca2021}. De uitleg mag dus geen overdreven zekerheid uitstralen en moet overeenkomen met de werkelijke beoordeling.
\end{description}
\subsection{Afbakening van de scope}
\label{subsec:afbakening-scope}
De doelstelling van de oplossing is om aan een e-mail een risicoscore te hangen die aangeeft hoe sterk de gegeven input op phishing wijst. Het formaat van die input, volgt de RFC~5322 zoals opgesteld door \textcite{Resnick2008}. 

Gelet op de eis van uitlegbaarheid, wordt gekozen voor een interpreteerbaar model. Zoals besproken in sectie~\ref{subsubsec:logistische-regressie} is logistische regressie een geschikte kandidaat, aangezien elk kenmerk een gewicht krijgt dat rechtstreeks aangeeft hoe het de risicoscore beïnvloedt. Verder kan er voor elke beslissing vooraf een uitleg worden gedefinieerd die dan in het uiteindelijke antwoord kan worden gecapteerd.

De kenmerken waarop het model getraind zal worden, worden berekend uit de datasets die verder worden besproken \ref{subsec:voorbereiding-data}. 

Wat buiten de scope van dit onderzoek valt en dus niet wordt opgenomen:
\begin{enumerate}
    \item \textbf{Bijlagen als detectiekenmerk.} Zoals aangehaald in sectie \ref{subsubsec:bijlagen} blijken bijlagekenmerken volgens \textcite{Mendes2025} slechts een beperkte bijdrage te hebben in het eindoordeel van het model. De geringe meerwaarde weegt niet op tegen de bijkomende complexiteit;
    \item \textbf{LLM's voor beoordeling en uitleg.} Op basis van de bevindingen van \textcite{Turpin2023} valt een LLM af als generator van uitleg: de uitleg moet waarheidsgetrouw zijn, en dat kan een LLM niet in elk geval garanderen. Ook \emph{indirect prompt injection} \autocite{Greshake2023} (sectie~\ref{subsubsec:machinelearning-detectie:nadelen}) vormt een risico voor de correcte werking van de oplossing, aangezien de inhoud van een phishingmail door een aanvaller wordt bepaald en dus ongewenste instructies voor de LLM kan bevatten;
    \item \textbf{Inhoudelijke (lexicale) kenmerken.} GenAI maakt foutloze en overtuigende teksten makkelijk op te stellen, waardoor taalkenmerken minder betrouwbaar worden dan structurele kenmerken. Verder zijn de beschikbare corpora overwegend Engelstalig \autocite{Pajola2025} (sectie~\ref{subsubsec:datasets}), waardoor een model dat op woordgebruik wordt getraind niet naar het Nederlands bruikbaar zou zijn.
    \item \textbf{Stilometrische kenmerken.} Deze kenmerken vergelijken het schrijfgedrag in een bericht met de gebruikelijke schrijfstijl van de afzender. Dat vereist voldoende eerdere berichten van dezelfde afzender \autocite{Stringhini2014} (sectie~\ref{subsubsec:inhoud-taalgebruik:stilometrische-kenmerken}), terwijl het algoritme telkens één afzonderlijk bericht beoordeelt.
    \item \textbf{Authenticatiekenmerken.} De resultaten van SPF-, DKIM- en DMARC-controles (sectie~\ref{subsubsec:Afzender en e-mailheaders:authenticatiekenmerken}) worden niet als kenmerk opgenomen. De legitieme e-mails uit het SpamAssassin Public Corpus dateren van vóór de wijdverspreide invoering van deze mechanismen en bevatten dus geen \texttt{Authentication-Results}-header, terwijl recente phishingmails die wel kunnen bevatten. Het model zou dan de verkeerde conclusie trekken dat de aanwezigheid van die header wijst op phishing.
\end{enumerate}
\subsection{Voorbereiding data}
\label{subsec:voorbereiding-data}

De in sectie~\ref{subsubsec:datasets} besproken corpora zijn niet geschikt voor dit onderzoek, omdat ze niet de vereiste e-mailheaders bevatten. Daarom wordt geen gebruikgemaakt van MeAJOR \autocite{Mendes2025}, Enron \autocite{klimt2004} of Nazario \autocite{Nazario2025}, maar van de volgende twee datasets:
\begin{enumerate}
    \item \textbf{Phishing\_pot.} Een verzameling phishingmails die actief onderhouden wordt door \textcite{PhishingPot};
    \item \textbf{SpamAssassin Public Corpus.} Een verzameling e-mailberichten uit de vroege jaren 2000, met zowel ongewenste als legitieme e-mails \footfullcite{SpamAssassinCorpus}. Voor dit onderzoek wordt echter uitsluitend gebruikmaakt van de partities \\ \emph{easy\_ham} en \emph{hard\_ham}. De eerste bevat legitieme e-mails die relatief eenvoudig te detecteren zijn. De tweede bevat legitieme e-mails die meer gelijkenissen vertonen met frauduleuze berichten.
\end{enumerate}

\paragraph{Beperkingen van de data.}
Het SpamAssassin-project raadt zelf af om hun legitieme e-mailset te gebruiken als die ouder is dan 18 maanden is, omdat legitieme e-mail met de tijd verandert \footfullcite{SpamAssassinCorpusPolicy}.
Het gebruik van deze data blijft echter verdedigbaar, omdat het doel een risicoscore is op basis van structurele kenmerken, zoals headerconsistentie, URL-opbouw of hyperlinks. Deze kenmerken evolueren naar verwachting trager dan woordgebruik.

\paragraph{Voorbereiding.} 
Elke e-mail wordt ingelezen en omgezet naar een verwerkbare vorm. Daarbij worden headers en body gedecodeerd waar nodig, zoals \emph{encoded-words} in headers (RFC~2047 \autocite{Moore1996}) en body-delen die als base64 zijn gecodeerd. 

Vervolgens wordt de dataset opgesplitst in een trainingsset van 70\% en een testset van 30\%. Hierbij blijft de verhouding tussen phishingmails en legitieme e-mails in beide sets behouden. Een vaste seed zorgt ervoor dat de opsplitsing bij herhaaldelijke uitvoering reproduceerbaar is.

Uit elk bericht worden vervolgens de kenmerken berekend. Deze kenmerken, samen met de indeling van de berichten in trainings- en testset, worden opgeslagen. Zo hoeven de berichten niet bij elk experiment opnieuw te worden verwerkt, en worden alle modelconfiguraties op dezelfde gegevens getraind.

De testset wordt pas gebruikt wanneer de parameters van het model op punt staan na voldoende trainingsiteraties.


\subsection{Kenmerkextractie}%
\label{sec:kenmerkextractie}

Uit elk bericht worden de kenmerken afgeleid die in sectie~\ref{subsec:afbakening-scope} binnen de scope vallen, voor zover ze rechtstreeks uit de datasets van sectie~\ref{subsec:voorbereiding-data} komen of eruit af te leiden zijn. Het gaat om vier families die besprokekn werden in sectie~\ref{subsec:kenmerken-detectie}. 

De kenmerken worden eerst allemaal berekend als kandidaat-kenmerken. Welke ervan in het uiteindelijke model blijven, wordt bepaald aan de hand van L1-regularisatie zoals aangehaald in sectie~\ref{subsubsec:machinelearning-detectie}.

\paragraph{Opbouw van de extractors.}
Elke extractor levert twee elementen: een numerieke waarde voor het model en het concrete bewijs in de e-mail waarop die waarde berust. Zo kan de uitleg later controleerbaar naar de e-mail zelf verwijzen, zonder dat daarvoor een apart proces nodig is.

\paragraph{Afzenderidentiteit.}
De kenmerken in tabel~\ref{tab:kenmerken-afzender} volgen uit sectie~\ref{subsubsec:Afzender en e-mailheaders}. De controle op domeinen die lijken op legitieme domeinen en de controle op merken in de weergavenaam (A7) hangen af van een referentielijst.

\input{inhoud/methodologie/tabellen/3.1.tabel-kenmerken-afzender.tex}

\paragraph{Transportkenmerken.}
De kenmerken in tabel~\ref{tab:kenmerken-transport} worden eveneens besproken in sectie~\ref{subsubsec:Afzender en e-mailheaders:transport}.

\input{inhoud/methodologie/tabellen/3.2.tabel-kenmerken-transport.tex}

Kenmerk B4 telt het aantal hops in de \texttt{Received}-keten. Kenmerk B5 meet het verschil tussen het \texttt{Date}-veld en het oudste \texttt{Received}-tijdstempel. Omdat \texttt{Date} door de afzender zelf wordt ingesteld en de \texttt{Received}-tijdstempels door de betrokken servers worden toegevoegd, kan een groot verschil wijzen op manipulatie. Kenmerk B6 controleert ten slotte of de tijdstempels van opeenvolgende hops chronologisch oplopen, aangezien een afwijkende volgorde kan wijzen op een toegevoegde of aangepaste \texttt{Received}-header.

\paragraph{Lexicale URL-kenmerken.}
De kenmerken in tabel~\ref{tab:kenmerken-url} volgen uit sectie~\ref{subsubsec:urls-hyperlinks}.

\input{inhoud/methodologie/tabellen/3.3.tabel-url.tex}

C9 steunt op de private sectie van de Public Suffix List, die onder meer gratis hostingplatformen bevat die ook door aanvallers worden misbruikt.

\paragraph{Misleidende hyperlinks.}
De kenmerken in tabel~\ref{tab:kenmerken-url} volgen uit sectie~\ref{subsubsec:urls-hyperlinks::misleidende-hyperlinks}.
Deze kenmerken zijn enkel van toepassing op mails waarbij de body een stuk HTML bevat [waarom?]

\input{inhoud/methodologie/tabellen/3.4.tabel-hyperlinks.tex}

\paragraph{Niet-bepaalbare kenmerken.}
Sommige kandidaat-kenmerken zijn niet voor elke e-mail te bepalen, bijvoorbeeld omdat het bericht geen HTML-deel of geen \texttt{Reply-To} bevat. In dat geval valt er geen kandidaat-kenmerk te berekenen, en dat wordt expliciet als \enquote{niet van toepassing} vastgelegd. Dit mag niet gelijkgesteld worden met een waarde van 0 of false, want dit betekent dat een controle werd uitgevoerd en geen afwijking vond. Hoe vaak een kandidaat-kenmerk niet bepaalbaar was, wordt apart bijgehouden voor de rapportering achteraf, maar niet als modelkenmerk gebruikt.
\subsection{Uitleg}
\label{subsec:uitleg}

De uitleg beantwoordt de vraag waarom het systeem een e-mail als risicovol beoordeelt en wordt in gewone taal aan de eindgebruiker voorgelegd. Zoals besproken in sectie~\ref{subsubsec:vormen-van-uitlegbaarheid} is hiervoor een lokale uitleg het meest geschikt, aangezien die de beoordeling van één individuele e-mail verduidelijkt \autocite{Guidotti2018}.

Bij het opstellen van de uitlegregels worden de volgende uitgangspunten gehanteerd, die requirement R2 tot en met R4 vertalen uit sectie~\ref{subsec:requirements}:
\begin{enumerate}
    \item \textbf{Suggestief.} De verantwoordelijkheid voor de interpretatie blijft bij de eindgebruiker. De uitlegregels mogen dus geen definitief oordeel over de e-mail uitbrengen, maar dienen de gebruiker met voorzichtige formuleringen zoals ``kan wijzen op'' te informeren;
    \item \textbf{Waarheidsgetrouw.} Elke reden komt rechtstreeks uit de berekening van de score en berust op kenmerken die in het bericht daadwerkelijk worden vastgesteld. Kenmerken die niet of onvoldoende aanwezig zijn, worden niet als reden opgenomen;
    \item \textbf{Controleerbaar.} Elke reden verwijst naar concreet bewijs uit het bericht. Dat bewijs wordt samen met de kenmerkwaarde door de extractors bijgehouden (sectie~\ref{sec:kenmerkextractie}), zodat de gebruiker elke reden zelf in de e-mail kan nagaan;
     \item \textbf{Toegankelijk.} De uitleg is geschreven in gewone taal en vermijdt jargon.
\end{enumerate}

\paragraph{Kanttekening bij uitleg in deze scriptie.}
De uitleg is een essentieel onderdeel van het algoritme, want ondersteuning van de gebruiker vraagt niet enkel om een score, maar ook om de redenen erachter (sectie~\ref{subsec:uitlegbaarheid-en-gebruikersondersteuning}). Dit onderzoek toont aan dat die redenen rechtstreeks uit de berekening van de score kunnen worden afgeleid, naar concreet bewijs in de e-mail verwijzen en in gewone taal kunnen worden weergegeven. De formulering wordt niet bij gebruikers getoetst en het effect op hun begrip, vertrouwen en beslissingen, inclusief \emph{overreliance} (sectie~\ref{subsubsec:uitleg-voor-gebruikers}), valt buiten de scope.
\subsection{Evaluatie}
\label{sec:evaluatie}

Om de bruikbaarheid van het algoritme te evalueren, wordt gebruikgemaakt van de testset die tijdens de voorbereiding van de data werd afgesplitst (sectie~\ref{subsec:voorbereiding-data}). De resultaten van het model worden gerapporteerd aan de hand van de metrieken uit sectie~\ref{subsubsec:evaluatiemetrieken}, met het false-positive-percentage (\emph{FPR}) als de belangrijkste. Deze metriek geeft immers aan hoe vaak het algoritme een legitieme e-mail ten onrechte als phishing aanduidt. Zoals besproken in sectie~\ref{subsubsec:uitleg-voor-gebruikers}, besteden gebruikers steeds minder aandacht aan waarschuwingen die ze herhaaldelijk te zien krijgen \autocite{Anderson2015}. Een hoge \emph{FPR} brengt dan ook het risico mee dat gebruikers aan de waarschuwingen van het algoritme wennen, waardoor dit onderzoek zijn doel mist.




%---------- Verwachte resultaten ----------------------------------------------
\section{Verwacht resultaat, conclusie}%
\label{sec:verwachte_resultaten}

Hier beschrijf je welke resultaten je verwacht. Als je metingen en simulaties uitvoert, kan je hier al mock-ups maken van de grafieken samen met de verwachte conclusies. Benoem zeker al je assen en de onderdelen van de grafiek die je gaat gebruiken. Dit zorgt ervoor dat je concreet weet welk soort data je moet verzamelen en hoe je die moet meten.

Wat heeft de doelgroep van je onderzoek aan het resultaat? Op welke manier zorgt jouw bachelorproef voor een meerwaarde?

Hier beschrijf je wat je verwacht uit je onderzoek, met de motivatie waarom. Het is \textbf{niet} erg indien uit je onderzoek andere resultaten en conclusies vloeien dan dat je hier beschrijft: het is dan juist interessant om te onderzoeken waarom jouw hypothesen niet overeenkomen met de resultaten.



\printbibliography[heading=bibintoc]

\end{document}


Hoe kan een interpreteerbare risicoscore voor phishingmails worden opgebouwd uit structurele e-mailkenmerken, zodat phishing met een lage false-positive-rate wordt gedetecteerd en elke score herleidbaar is tot controleerbaar bewijs in de e-mail?